\section{Jacobsthal Sums}
Building further on the theme of character sums twisted by the quadratic character, we now introduce a classical family of sums first studied by Ernst Jacobsthal in 1907 in connection with the representation of primes $p\equiv 1\bmod 4$ as sums of two squares. Unlike Gaussian and Jacobi sums, which are defined purely multiplicatively, Jacobsthal sums mix an additive shift with the quadratic character $\eta$ of $\bbF_q$, which makes them particularly well suited to problems involving quadratic residues and the arithmetic of $\bbF_q^*$.
\begin{Definition}{Jacobsthal Sums}{}
  For $\alpha\in\bbF_q^*$ where $q$ is odd and for any $n\in\bbN$ the \emph{Jacobsthal Sum} is defined as $$H_n(\alpha)=\sum_{\gm\in\bbF_q}\eta(\gm^{n+1}+\alpha\cdot \gm)=\sum_{\gm\in\bbF_q}\eta(\gm)\cdot \eta(\gm^n+\alpha)$$
\end{Definition}